\answer{ex:20:1}{$t_{\text{w}}=\tau_r\ln\frac{1-L}{1-H}=16.8$~h、$t_{\text{s}}=\tau_d\ln\frac HL=5.8$~h、周期$22.5$~h。発振器を$24$~hに引き込むのは閾値の概日的な揺れ、すなわち位相同期~\eqref{eq:pll}である。}
\answer{ex:20:2}{(a)~$L=0.111$。$L=0.092$なら睡眠は$9.6$~h続く。(b)~$18\%$ではなく$22\%$（$1/0.82=1.22$）。}
\answer{ex:20:3}{$H=\frac{p-1}{p-1/q}=0.623$、$L=H/q=0.093$、ただし$p=e^{16/\tau_r}$、$q=e^{8/\tau_d}$。$4$~hで起こされると位相は永久にずれる（入眠が$7.2$~h早まる）。閾値が揺れるなら位相は2--3日で戻る。}
\answer{ex:20:4}{(a)~$r_\pm=\frac12\bigl(1\pm\sqrt{1-4k/w}\bigr)=0.947,\ 0.053$。$a_{\text{up}}=3.51$、$a_{\text{down}}=0.49$。(b)~活動$\approx0.33$~s、静止$\approx0.59$~s、周期$\approx0.93$~s、活動の割合$0.36$。}
\answer{ex:20:5}{$a_{\text{up}}/r_+<b<a_{\text{down}}/r_-$：$3.71<b<9.20$。\nt{ACET}は$b$を$3.71$未満に下げる。交点は上の枝に移り、活動状態（$r\approx1$）が安定になる。}
\answer{ex:20:6}{$D=24$、$32$、$40$、$48$、$64$で$4.9$、$6.8$、$8.8$、$5.5$、$8.9$~h。長さを決めるのは借りではなく、入眠の概日位相である。}
\answer{ex:20:7}{Bのあと、Aでの誤差は平均$0.51$（$20$セットで$0.19$から$1.08$。応答ゼロなら$1$）。混ぜて学習すると、どちらの誤差も$10^{-4}$未満。}
\answer{ex:20:8}{自由課題。一様なスケーリングは重みの比を保つが、閾値ニューロンの応答を保つのは閾値も同じ比率でスケーリングされる場合に限る。インターリーブしたリプレイは比を変える。大きな接触が変わらないことは、純粋に一様なスケーリングと矛盾する。}
